\section{支撑材料文件清单}\label{app:materials}
表\ref{tab:catalog}列出当前支撑材料文件及用途。工作区根目录的“答案”文件夹集中保存题目要求填写的四个结果工作簿；其余材料位于“A题/论文提升\_20260913\_v4/”工程目录，表中其余路径均相对于该工程目录。附录B给出七个脚本、五个辅助模块及依赖清单。

\begin{table}[htbp]\centering\small
\caption{电子支撑材料目录}\label{tab:catalog}
\begin{tabularx}{\linewidth}{lX}
\toprule
文件或目录 & 内容\\
\midrule
答案 & \texttt{result1.xlsx}、\texttt{result2.xlsx}、\texttt{result3.xlsx}、\texttt{result4.xlsx}，分别对应问题一至四的已填写结果，与\texttt{results/}中的同名文件一致\\
\texttt{solve.py} & 物性、网格、通量、雅可比、积分、事件与小算例检验\\
\texttt{run\_all.py} & 网格和时间复验、结果表、场数组及情景分析\\
\texttt{make\_figures.py} & 从真实结果生成14组矢量和位图\\
\texttt{verify\_moving\_domain.py} & 移动域解析模态、时空加密、几何因子对照与守恒检验\\
\texttt{export\_result2\_full.py} & 全过程逐秒工作簿的分块导出和完整结构检查\\
\texttt{reproduce.py} & 完整计算、验证和导出的运行入口\\
\texttt{test\_sync.py} & 采样统计与600 s导出的短程检查\\
\texttt{utils/} & 全量读表、绘图风格、图表自检、导出与环境记录\\
\texttt{requirements.txt} & 原计算使用的Python包版本\\
\texttt{results/} & 四个结果工作簿的原始输出、六张题表、全过程时空输出、过程数组、域外空值规则、诊断与参数摘要\\
\texttt{figures/} & 原始输入、运行过程和结果图的PDF、SVG、PNG\\
新增示意图/ & 三张框架、流程与随体坐标图的可编辑源文件和矢量插入版本\\
完整论文-LaTeX/ & 主文稿、章节、表格、图形及完整代码副本\\
AI工具使用详情.pdf & 关键提示、采纳范围、工具说明和人工复核提示\\
AI工具说明-LaTeX/ & AI工具使用说明的LaTeX源码与配图\\
\bottomrule
\end{tabularx}
\end{table}

计算环境为Python 3.10.12，包版本见\texttt{requirements.txt}。准备依赖及文泉驿正黑字体\texttt{wqy-zenhei.ttc}后，从包含原始附件的项目根目录执行：
\begin{Verbatim}[fontsize=\small,breaklines,breakanywhere]
OPENBLAS_NUM_THREADS=1 OMP_NUM_THREADS=1 python3 -B reproduce.py
\end{Verbatim}

完整入口在\texttt{reproduced\_*}目录依次计算解析基准、四问结果、空间与时间精度比较、工作簿、情景对照和图形，保留原有结果。需要短程检查或单独运行移动域算例时，可分别使用：
\begin{Verbatim}[fontsize=\small,breaklines,breakanywhere]
python3 -B reproduce.py --smoke
python3 -B verify_moving_domain.py
python3 -B verify_moving_domain.py --mutant-only
\end{Verbatim}

\section{完整可运行源代码}

\subsection{核心求解与诊断}
\VerbatimInput[fontsize=\scriptsize,breaklines,breakanywhere,breaksymbolleft={},breaksymbolright={},frame=single,framerule=0.4pt,framesep=2mm]{code/solve.py}
\subsection{全流程计算与结果导出}
\VerbatimInput[fontsize=\scriptsize,breaklines,breakanywhere,breaksymbolleft={},breaksymbolright={},frame=single,framerule=0.4pt,framesep=2mm]{code/run_all.py}
\subsection{真实结果可视化}
\VerbatimInput[fontsize=\scriptsize,breaklines,breakanywhere,breaksymbolleft={},breaksymbolright={},frame=single,framerule=0.4pt,framesep=2mm]{code/make_figures.py}
\subsection{移动域解析基准与几何因子对照}
\VerbatimInput[fontsize=\scriptsize,breaklines,breakanywhere,breaksymbolleft={},breaksymbolright={},frame=single,framerule=0.4pt,framesep=2mm]{code/verify_moving_domain.py}
\subsection{全过程逐秒工作簿导出与核对}
\VerbatimInput[fontsize=\scriptsize,breaklines,breakanywhere,breaksymbolleft={},breaksymbolright={},frame=single,framerule=0.4pt,framesep=2mm]{code/export_result2_full.py}
\subsection{完整计算入口}
\VerbatimInput[fontsize=\scriptsize,breaklines,breakanywhere,breaksymbolleft={},breaksymbolright={},frame=single,framerule=0.4pt,framesep=2mm]{code/reproduce.py}
\subsection{采样与导出检查}
\VerbatimInput[fontsize=\scriptsize,breaklines,breakanywhere,breaksymbolleft={},breaksymbolright={},frame=single,framerule=0.4pt,framesep=2mm]{code/test_sync.py}
\subsection{完整附件读取}
\VerbatimInput[fontsize=\scriptsize,breaklines,breakanywhere,breaksymbolleft={},breaksymbolright={},frame=single,framerule=0.4pt,framesep=2mm]{code/utils/read_rows.py}
\subsection{图表结构与布局检查}
\VerbatimInput[fontsize=\scriptsize,breaklines,breakanywhere,breaksymbolleft={},breaksymbolright={},frame=single,framerule=0.4pt,framesep=2mm]{code/utils/plot_style.py}
\subsection{统一科研绘图风格}
\VerbatimInput[fontsize=\scriptsize,breaklines,breakanywhere,breaksymbolleft={},breaksymbolright={},frame=single,framerule=0.4pt,framesep=2mm]{code/utils/setup_style.py}
\subsection{多格式图表导出}
\VerbatimInput[fontsize=\scriptsize,breaklines,breakanywhere,breaksymbolleft={},breaksymbolright={},frame=single,framerule=0.4pt,framesep=2mm]{code/utils/export_figure.py}
\subsection{输入与运行环境记录}
\VerbatimInput[fontsize=\scriptsize,breaklines,breakanywhere,breaksymbolleft={},breaksymbolright={},frame=single,framerule=0.4pt,framesep=2mm]{code/utils/repro_manifest.py}
\subsection{Python依赖版本}
\VerbatimInput[fontsize=\small,breaklines,frame=single,framerule=0.4pt,framesep=2mm]{code/requirements.txt}
